\documentclass[11pt]{article}
\usepackage[utf8]{inputenc}
\usepackage[T1]{fontenc}
\usepackage[italian]{babel}
\usepackage{geometry}
\geometry{a4paper, margin=1in}
\usepackage{xcolor}
\usepackage{listings}
\usepackage{hyperref}

\usepackage{tcolorbox}
\tcbuselibrary{breakable}
\begin{document}

\begin{tcolorbox}[colback=blue!5!white,colframe=blue!75!black,title=\textbf{DOMANDA}, breakable]
Esercizi su TCP\\
11) Lanciare due volte il server usando due terminali. Cosa si osserva? Funzionano entrambi?
\end{tcolorbox}

\begin{tcolorbox}[colback=green!5!white,colframe=green!75!black,title=\textbf{RISPOSTA}, breakable]
Essendo ora passati agli esercizi su \texttt{TCP}, il comportamento del programma è differente rispetto al precedente caso \texttt{UDP}.

Lanciando due volte il server \texttt{TCP} (presente nella cartella es12) su due terminali, si osserva quanto segue:

\begin{itemize}
    \item \textbf{Il primo server funziona regolarmente:} si avvia, esegue il \texttt{bind()} e si mette in ascolto delle connessioni (\texttt{listen()}), stampando il messaggio di attesa a video.
    \item Il secondo server NON funziona e termina immediatamente.
\end{itemize}

\textbf{Perché succede questo?} A differenza di quanto visto con \texttt{UDP}, nel protocollo \texttt{TCP} l'opzione del socket \texttt{SO\_REUSEADDR} ha un significato più stringente: serve esclusivamente a consentire il riutilizzo di una porta che si trova in stato di "TIME\_WAIT" (cioè una porta utilizzata da una connessione chiusa di recente dal sistema operativo). Tuttavia, \texttt{SO\_REUSEADDR} non consente a due programmi di mettersi contemporaneamente in stato di ascolto attivo (\texttt{listen}) sulla medesima porta e \texttt{IP}. Di conseguenza, il tentativo di eseguire la \texttt{bind()} da parte del secondo server andrà incontro a un errore (del tipo Address already in use) e, come previsto dalla gestione degli errori all'interno della funzione \texttt{createTCPServer} nel file \texttt{network.c}, causerà la chiusura immediata (terminazione silente tramite \texttt{exit()}) del secondo programma.
\end{tcolorbox}

\end{document}
